%==========================================================
\section{Avaliação do protótipo}\label{sec:avaliacao}

O protótipo foi avaliado de duas formas complementares, por um elemento da equipa que não o construiu: uma avaliação heurística, com as dez heurísticas de Nielsen~\cite{nielsen1990,nielsen1994}, e um percurso cognitivo~\cite{wharton1994}.

\begin{table}[!htbp]
\centering\small
\begin{tabularx}{\textwidth}{@{}p{4cm}L@{}}
\toprule
Protótipo construído por & \textit{a preencher} \\
Avaliação heurística feita por & \textit{a preencher} \\
Percurso cognitivo feito por & \textit{a preencher} \\
Data & \textit{a preencher} \\
\bottomrule
\end{tabularx}
\caption{Quem construiu e quem avaliou.}
\label{tab:avaliadores}
\end{table}

%----------------------------------------------------------
\subsection{Avaliação heurística}\label{sec:heuristica}

Cada problema encontrado é associado a uma heurística, a um ecrã do protótipo e a uma gravidade na escala de Nielsen: 0 -- não é problema; 1 -- cosmético; 2 -- menor; 3 -- maior; 4 -- catastrófico.

\begin{table}[!htbp]
\centering\small
\begin{tabularx}{\textwidth}{@{}p{0.6cm}p{5.4cm}Lc@{}}
\toprule
\textbf{H} & \textbf{Heurística} & \textbf{Problemas encontrados} & \textbf{N.º} \\
\midrule
H1 & Visibilidade do estado do sistema & & \\
H2 & Correspondência entre o sistema e o mundo real & & \\
H3 & Controlo e liberdade do utilizador & & \\
H4 & Consistência e normas & & \\
H5 & Prevenção de erros & & \\
H6 & Reconhecer em vez de recordar & & \\
H7 & Flexibilidade e eficiência de uso & & \\
H8 & Estética e design minimalista & & \\
H9 & Ajudar a reconhecer, diagnosticar e recuperar de erros & & \\
H10 & Ajuda e documentação & & \\
\bottomrule
\end{tabularx}
\caption{Resumo da avaliação heurística por heurística.}
\label{tab:heuristica-resumo}
\end{table}

\begin{table}[!htbp]
\centering\small
\begin{tabularx}{\textwidth}{@{}p{0.9cm}p{0.8cm}p{1cm}Lc@{}}
\toprule
\textbf{ID} & \textbf{H} & \textbf{Ecrã} & \textbf{Problema e recomendação} & \textbf{Grav.} \\
\midrule
HE01 & & & \textit{a preencher} & \\
HE02 & & & & \\
\bottomrule
\end{tabularx}
\caption{Problemas encontrados na avaliação heurística.}
\label{tab:heuristica}
\end{table}

\pendente{a definir (quem não construiu o protótipo)}{%
Preencher as Tabelas~\ref{tab:avaliadores}, \ref{tab:heuristica-resumo} e~\ref{tab:heuristica}. Se a IA gerar uma primeira avaliação a partir das imagens do protótipo, guardá-la e comparar a sua precisão e cobertura com a avaliação humana (Apêndice~\ref{anx:ia}).}

%----------------------------------------------------------
\subsection{Percurso cognitivo}\label{sec:percurso}

O percurso cognitivo segue, passo a passo, um utilizador novo a cumprir uma tarefa concreta. A tarefa e o utilizador devem vir dos dados de campo, por exemplo: \textit{«U0 quer saber quanto custa, ao todo e por pessoa, uma viagem de 4 noites para três amigos dentro de 300\,€ por pessoa, e partilhá-la com o grupo.»}

\begin{table}[!htbp]
\centering\small
\begin{tabularx}{\textwidth}{@{}p{3.4cm}L@{}}
\toprule
Utilizador & \textit{a preencher (perfil baseado em U0 ou noutro participante)} \\
Tarefa & \textit{a preencher} \\
Sequência de ações correta & \textit{a preencher} \\
\bottomrule
\end{tabularx}
\caption{Definição do percurso cognitivo.}
\label{tab:percurso-def}
\end{table}

Em cada passo respondem-se às quatro perguntas do método~\cite{wharton1994}: (Q1) o utilizador vai tentar alcançar o efeito certo? (Q2) vai reparar que a ação certa está disponível? (Q3) vai associar a ação certa ao efeito que quer? (Q4) depois da ação, vai perceber que avançou na tarefa?

\begin{table}[!htbp]
\centering\small
\begin{tabularx}{\textwidth}{@{}p{0.6cm}p{3cm}ccccL@{}}
\toprule
\textbf{N.º} & \textbf{Ação} & \textbf{Q1} & \textbf{Q2} & \textbf{Q3} & \textbf{Q4} & \textbf{Problema (se algum «não»)} \\
\midrule
1 & \textit{a preencher} & & & & & \\
2 & & & & & & \\
3 & & & & & & \\
\bottomrule
\end{tabularx}
\caption{Percurso cognitivo, passo a passo.}
\label{tab:percurso}
\end{table}

\pendente{a definir (quem não construiu o protótipo)}{%
Preencher as Tabelas~\ref{tab:percurso-def} e~\ref{tab:percurso}. Se a IA redigir o guião da tarefa, guardar a versão original e a revista (Apêndice~\ref{anx:ia}).}
